%%=============================================================================
%% De inhoud van je projectvoorstel
%%=============================================================================
%%
%% Zeven rubrieken, samen hoogstens twee bladzijden. Houd het kort: per
%% rubriek volstaan één tot drie zinnen. Een voorstel is geen rapport - het
%% moet je begeleider in vijf minuten laten beslissen of dit een geschikt
%% graduaatsproject is.
%%
%% Stem je voorstel af met je mentor VOORDAT je het indient. Je mentor
%% bevestigt dat dit is wat het bedrijf bedoelt en dat de tien werkdagen
%% haalbaar zijn.

%---------- 1. Context en probleemstelling ------------------------------------

\section{Context en probleemstelling}%
\label{vst:context}

Ik loop stage bij CMC NV. er is één software team die de software onderhoudt van een aantal producten. Deze producten worden gebruikt in de bedrijfswereld, en als er een probleem is met een product moet er een RMA (return merchandise authorization) aangemaakt worden
om dit te kunnen verwezelijken. Het huidige process is heel verouderd en vergt heel veel manueel werk, also de server waar dit huidige process loopt is heel verouderd, waardoor deze op elk moment offline worden gehaald.
Hier kruipt zeer veel werk in, dus er moet een site worden gemaakt waar alles veel meer geautomiseerd en klantvriendelijk is.

%---------- 2. Doelstelling ---------------------------------------------------

\section{Doelstelling}%
\label{vst:doelstelling}

Op het einde van het project bestaat een website waar men als klant zijnde gemakkelijk een RMA indient. De validatie, bijhouding van de status, etc... zal veel meer geautomiseerd zijn dan het nu is, en er zal veel minder manueel bijgehouden worden.

%---------- 3. Technologieën --------------------------------------------------

\section{Technologieën}%
\label{vst:technologieen}

Er zal gebruik gemaakt worden van de programeertaal C\#, voor de frameworks zal ik gebruik maken van Blazor WASM en EF Core, voor databanken zal ik gebruik maken van Microsoft SQL Server, verder zal ik ook gebruik maken van Github, Keycloak en Uvion Uhost.

Ik veronderstel AI te gebruiken, en dit wordt ook aangemoedigd door mijn stagebedrijf. Ik veronderstel dit vooral in te zetten voor syntax op te zoeken waar ik nog niet vertrouwd mee ben, alsook vooral voor bugfixing.

%---------- 4. Afbakening -----------------------------------------------------

\section{Afbakening}%
\label{vst:afbakening}

Mijn dagelijks stagewerk is meehelpen aan de ontwikkeling van een reporting engine voor rapporten aan te maken via de verschillende applicaties van het bedrijf.

Het project dat ik zelfstanding zou uitwerken is een volledige herwerking van het huidige RMA process, analyse + ontwerp van hoe de website er zou uitzien + volledige nieuwe codebase uitschrijven voor nieuwe website.

%---------- 5. Verwachte meerwaarde -------------------------------------------

\section{Verwachte meerwaarde}%
\label{vst:meerwaarde}

Het RMA process zal veel vlotter en veel geautomatiseerder verlopen. Doordat da status van alle RMA's niet meer manueel moet bijgehouden worden, zal dit voor veel tijdswinst zorgen en tot minder fouten lijden.

%---------- 6. Planning en werkbelasting --------------------------------------

\section{Planning en werkbelasting}%
\label{vst:planning}

% Het graduaatsproject is ongeveer tien werkdagen, samen ongeveer 100 uren.
% Verdeel die over de fasen van je project en geef per fase aan wat je oplevert.
% Vermeld ook of je die dagen gefocust opneemt of gespreid over de stageperiode,
% en dat je dit met je mentor afgesproken hebt.

\begin{table}[htbp]
  \centering
  \begin{tabular}{llr}
    \toprule
    \textbf{Fase} & \textbf{Wat lever je op} & \textbf{Uren} \\
    \midrule
    Analyse en vereisten & overzicht van de vereisten   & 5  \\
    Technologiekeuze     & verantwoorde keuze           & 5  \\
    Ontwerp              & architectuur en gegevensmodel& 20  \\
    Realisatie           & werkende oplossing           & 50  \\
    Testen en validatie  & testverslag                  & 10  \\
    Rapport, poster en demo & de op te leveren documenten  & 10  \\
    \midrule
    \textbf{Totaal}           &                              & \textbf{100} \\
    \bottomrule
  \end{tabular}
  \caption[Planning en werkbelasting]{\label{tab:vst-planning}%
    }
\end{table}

%---------- 7. Risico's -------------------------------------------------------

\section{Risico's}%
\label{vst:risicos}

1. Blazor is een nieuwe technologie waar ik nog niet vertrouwd mee ben: Meer tijd plannen om documentatie te doornemen. 
2. Problemen met hosten via Uvion. Een bepaalde zekerheid inbouwen in de planning waardoor problemen met hosting tijdig kunnen worden ontdekt en opgelost.
